\chapter{Resultados y discusión comparativa}
\label{ch:resultados-discusion}

\section{Introducción del capítulo}
Este capítulo reúne los hallazgos principales del estudio comparativo y los interpreta a la luz
del problema físico que se desea resolver. El foco no está solo en identificar qué algoritmo
produce el menor valor de costo, sino en entender cómo esa sintonización modifica la dinámica
de frecuencia, el amortiguamiento del sistema y el esfuerzo que debe asumir el actuador
hidráulico.

\section{Objetivo del capítulo}
El objetivo es presentar una lectura técnica integrada de los resultados obtenidos en los
escenarios de validación. Para ello se comparan tendencias de respuesta temporal, métricas
numéricas y observaciones de robustez, resaltando la forma en que los algoritmos de
optimización alteran el compromiso entre rapidez y estabilidad.

\section{Resultados por algoritmo}
De manera general, las sintonizaciones metaheurísticas muestran una mejora visible frente a la
línea base obtenida con Ziegler--Nichols. En particular, el enfoque basado en PSO tiende a
reducir el ITAE con mayor consistencia y a ofrecer trayectorias de frecuencia más limpias en los
escenarios de perturbación. GWO y ABC aportan resultados competitivos en distintas pruebas,
mientras que Eagle Strategy ofrece una referencia interesante cuando se privilegia una búsqueda
híbrida entre exploración amplia y refinamiento local.

\section{Análisis comparativo}
Las diferencias entre algoritmos no deben interpretarse únicamente como una carrera por el
menor valor numérico. En algunos casos, una sintonización más agresiva reduce el tiempo de
asentamiento, pero incrementa el esfuerzo de control o acerca la planta a zonas indeseables de
operación. La ventaja del PSO observada en el borrador de referencia se explica porque logra
incrementar la capacidad de corrección sin perder amortiguamiento, en buena medida gracias al
papel de la acción derivativa en la solución encontrada.

\section{Ventajas y desventajas observadas}
La principal ventaja de los métodos metaheurísticos es que adaptan la sintonización a la planta
realmente simulada y no a una representación crítica simplificada. Como contraparte, exigen un
presupuesto computacional mayor y una definición cuidadosa de hiperparámetros. Entre los
algoritmos comparados, PSO destaca por su sencillez de implementación; GWO por su liderazgo
múltiple; ABC por su mecanismo de reinicio de soluciones; y Eagle Strategy por su lectura
híbrida del proceso de búsqueda.

\section{Interpretación ingenieril}
La lectura física de los resultados es tan importante como la lectura matemática. Menores
oscilaciones implican menos movimientos bruscos del servomecanismo, menor fatiga mecánica
y una operación más amigable con la turbina Michell--Banki. Por ello, mejorar sobreimpulso y
tiempo de asentamiento no solo eleva la calidad de la energía entregada a la carga, sino que
también protege el hardware existente y reduce la necesidad de soluciones externas costosas,
como almacenamiento adicional o compensación inercial artificial.

\section{Cierre del capítulo}
Los resultados respaldan la conveniencia de abordar la sintonización PID desde una lógica de
optimización y no únicamente mediante reglas históricas. El capítulo final resume los hallazgos,
define el alcance real del estudio y propone extensiones para fortalecer su validación futura.
